% Generated from evidence/benchmark_summary.json; all times are milliseconds.
\begin{table}[htbp]
\centering
\small
\setlength{\tabcolsep}{3.8pt}
\begin{tabular}{r rrr rrr}
\toprule
& \multicolumn{3}{c}{Full construction and enumeration}
& \multicolumn{3}{c}{Enumeration on the same kernel} \\
\cmidrule(lr){2-4}\cmidrule(lr){5-7}
$n$ & Incumbent & Planar & Paired ratio & Incumbent & Planar & Paired ratio \\
\midrule
1 & 4.238 & 2.720 & 1.47 & 3.002 & 2.363 & 1.27 \\
2 & 8.550 & 3.057 & 2.83 & 8.224 & 2.898 & 2.84 \\
4 & 46.490 & 5.209 & 8.90 & 47.096 & 4.367 & 10.79 \\
8 & 537.411 & 8.730 & 61.64 & 530.946 & 6.916 & 73.80 \\
16 & 8\,572.669 & 19.861 & 482.82 & 8\,749.966 & 11.624 & 758.84 \\
32 & $40\,000$ cap & 48.458 & -- & $40\,000$ cap & 25.997 & -- \\
\bottomrule
\end{tabular}
\caption{Complete seven-ray enumeration for the double-capped Fibonacci family.
Times are medians of three alternating-order paired rounds, in milliseconds;
one warmup per arm is retained in the raw data and excluded from these medians.
Paired ratios are medians of within-round incumbent/planar time ratios,
not ratios of the displayed medians. For example, at $n=16$ the full paired
ratio is $482.82$, whereas the ratio of medians is $431.64$.
The two post-kernel arms receive exactly the same immutable incumbent-built kernel.
At $n=32$, all eight incumbent warmup/round runs reach the cooperative
$40\,000$\,ms cap after emitting seven rays but before proving enumeration
complete; the cap entries are allowances, not completed-run medians,
and no speed ratio is reported.}
\label{tab:planar-enumeration-benchmarks}
\end{table}

\begin{table}[htbp]
\centering
\small
\begin{tabular}{r rrrr}
\toprule
$n$ & Prior envelope (ms) & Planar (ms) & Paired prior/planar & Ratio of medians \\
\midrule
8 & 2.193 & 2.415 & 0.92 & 0.91 \\
32 & 8.534 & 8.741 & 0.97 & 0.98 \\
\bottomrule
\end{tabular}
\caption{Control against the earlier optimized nullity-two envelope enumerator,
loaded unmodified from the preceding research package. Both arms use the same
incumbent-built kernel and return the same three formula rays. Times are
medians of three paired rounds; values below one in either ratio column
indicate a slower planar implementation. These controls show no measured
advantage over the already optimized nullity-two method.}
\label{tab:planar-prior-envelope-control}
\end{table}

\begin{table}[htbp]
\centering
\small
\setlength{\tabcolsep}{4pt}
\begin{tabular}{r rrr rr}
\toprule
& \multicolumn{3}{c}{Complete construction (ms)}
& \multicolumn{2}{c}{Paired incumbent/sparse ratios} \\
\cmidrule(lr){2-4}\cmidrule(lr){5-6}
$k$ & Incumbent & Sparse once & Sparse reused & Reused source & Both cached \\
\midrule
1 & 72.133 & 81.463 & 8.004 & 8.84 & 1.88 \\
8 & 44.435 & 51.755 & 5.638 & 7.88 & 2.06 \\
32 & 61.126 & 64.297 & 12.041 & 4.74 & 2.02 \\
\bottomrule
\end{tabular}
\caption{Separate sparse-constructor measurements on supplied sparse sectors
of $t=1024$ layered solid tori. Times are medians of five paired rounds and
include unchanged exact rational elimination. The reused-source column
excludes its one-time preparation, while both ordinary single-call
constructors include preparation. The last column uses an additional control
where both constructors receive already validated geometry; it isolates
sparse kernel construction from validation reuse. Ratios are medians of
within-round ratios. The single-call sparse constructor is slower on these
cases. These are construction measurements, separate from complete ray
enumeration and from whole-knot recognition. All nine constructor cases
and every raw sample are retained in the data files.}
\label{tab:planar-sparse-source-control}
\end{table}
